\providecommand{\SysName}{CARE}

\section{Discussion}
\label{sec:discussion}

\noindent\textbf{Scope of \SysName{}.}
\SysName{} is designed for security-aware local refactoring of C/C++ code, not
for arbitrary program repair or large semantic rewrites. Its strongest use case
is a setting where an existing codebase already builds, the maintainer can
provide project-level validation commands, and the desired transformation is
localized around cleanup, resource lifecycle, validation checks, error-handling
paths, or other security-sensitive code patterns. In this scope, the LLM is not
treated as an autonomous developer. It proposes small edits under an explicit
safety contract, and the validator decides whether the candidate is acceptable.

\noindent\textbf{Limitations.}
The current prototype is conservative and incomplete in several ways. Its
static analyses are lightweight and do not prove full C/C++ semantics, deep
aliasing, concurrency behavior, or cross-translation-unit ownership contracts.
When clang or compile databases are unavailable, \SysName{} falls back to a
brace-aware parser, which improves coverage but reduces precision. The detector
front end can still produce noisy findings, especially for resource-like
patterns in tests, generated code, examples, or complex ownership idioms. The
evaluation therefore emphasizes validated patch yield rather than raw detector
counts: a finding becomes useful only when a generated patch survives the full
validation pipeline.

\noindent\textbf{Interpreting validation.}
Passing \SysName{} validation is evidence of safety, not a formal proof of
semantic equivalence. Build and test success depend on the target project's
available test suite, and the resource, semantic, and security checkers are
designed to catch common regressions rather than all possible regressions.
However, the ablation results show that these lightweight checkers are
practically important: without resource-aware validation, many unsafe candidates
would be accepted. Thus, \SysName{} should be viewed as a fail-closed
refactoring assistant whose accepted patches still deserve maintainer review,
but whose validation pipeline filters out many candidates that a direct LLM
workflow would surface.

\noindent\textbf{Future directions.}
The main engineering opportunities are improving detector precision, expanding
semantic validation, and making patch materialization more robust. Detector
precision can be improved by stronger ownership summaries, better path
relevance filters, and more clang-backed confirmation for high-impact findings.
Validation can be strengthened with generated regression tests, sanitizer-based
execution, symbolic or bounded checks for small functions, and richer
side-effect summaries. Finally, replacing raw unified-diff generation with a
more structured edit application layer could reduce patch-application failures
and increase the fraction of LLM proposals that reach meaningful validation.
